% !TEX program = XeLaTeX
\documentclass[UTF8,a4paper,12pt]{ctexart}

% HITSZ recommendation letter (Shenzhen campus).
% Adapted from the HIT LoR layout for Shenzhen University Town
% https://www.overleaf.com/latex/templates/hit-letter-of-recommendation-template/cchwdhcrhjtv

\usepackage{geometry}
\geometry{left=3.17cm,right=3.17cm,top=3.4cm,bottom=2.6cm}
\usepackage{xcolor}
\usepackage{tikz}
\usetikzlibrary{calc}
\usepackage{setspace}
\usepackage{hyperref}
\usepackage{parskip}

\definecolor{hitblue}{cmyk}{1,0.80,0,0}
\definecolor{hitlightblue}{cmyk}{0.35,0.10,0,0}

\pagestyle{empty}
\setstretch{1.28}
\hypersetup{colorlinks=true,urlcolor=hitblue}

\newcommand{\hitszletterhead}{%
  \begin{tikzpicture}[remember picture,overlay]
    \fill[hitblue] (current page.north west) rectangle ($(current page.north east)+(0,-2.35)$);
    \node[anchor=west,text=white,font=\bfseries\Large] at ($(current page.north west)+(1.15,-0.85)$)
      {哈尔滨工业大学（深圳）};
    \node[anchor=west,text=white,font=\small] at ($(current page.north west)+(1.15,-1.45)$)
      {Harbin Institute of Technology, Shenzhen};
    \node[anchor=east,text=white,font=\small\bfseries] at ($(current page.north east)+(-1.15,-0.85)$)
      {推荐信 / Recommendation};
    \node[anchor=east,text=white!90,font=\scriptsize] at ($(current page.north east)+(-1.15,-1.45)$)
      {计算机科学与技术学院};
    \fill[hitblue] ($(current page.south west)$) rectangle ($(current page.south east)+(0,0.28)$);
    \fill[hitlightblue] ($(current page.south west)+(0,0.28)$) rectangle ($(current page.south east)+(0,0.40)$);
    \node[anchor=south,text=hitblue,font=\scriptsize] at ($(current page.south)+(0,0.55)$)
      {深圳市南山区深圳大学城哈工大校区 518055 \quad Tel: +86-755-2603-3483};
  \end{tikzpicture}%
}

\begin{document}
\hitszletterhead
\vspace*{0.35cm}

\begin{flushright}
2026年3月18日
\end{flushright}

尊敬的招生委员会：

我是哈尔滨工业大学（深圳）计算机科学与技术学院教授王雪，现担任科学机器学习实验室负责人。我怀着十分肯定的态度，推荐陈远同学申请贵校博士研究生。自2024年春季起，陈远作为我的硕士研究生，全程参与国家自然科学基金面上项目“科学代理模型的残差校准与可信外推”。

陈远的研究工作把共形推断与可微物理代理结合起来。他在较小的标注校准集上估计单调运输映射，使 Burgers 方程、Darcy 流与反应--扩散系统上的区间覆盖率回到名义水平，并给出了当代理模型错过分岔时的失效检测方法。相关结果已整理为一篇中英双语预印本，陈远为第一作者。

在实验室日常工作中，陈远负责维护可复现实验流水线，并为两位本科生开放了校准模块的接口文档。他表达严谨，从不夸大覆盖率数字，这在数据驱动科学计算方向尤为难得。

综合学术潜力、工程能力与科研态度，我将陈远列为本实验室近五年硕士研究生中的前 5\%。我毫无保留地推荐他。如需补充材料，请随时与我联系。

\vspace{1.0em}
此致\\
敬礼

\vspace{1.3em}
\textbf{王雪 教授}\\
哈尔滨工业大学（深圳）计算机科学与技术学院\\
科学机器学习实验室\\
\href{mailto:wangxue@hit.edu.cn}{wangxue@hit.edu.cn} \quad +86-755-2603-3600

\vspace{1.4em}
\noindent\rule{\textwidth}{0.4pt}

\noindent\textbf{English summary.} I recommend Mr.\,Yuan Chen without reservation for doctoral study. He led a residual-score calibration project for scientific emulators, restored nominal interval coverage on three PDE families, and maintains a reproducible experimental pipeline. I rank him in the top 5\% of M.S.\ students I have supervised in the last five years.

\end{document}
